\documentclass[10pt,a4paper]{amsart}
\usepackage[margin=19mm]{geometry}
\usepackage{amsmath,amssymb,mathtools}
\usepackage[scr=rsfs,cal=euler]{mathalfa}
\usepackage{xfrac}
\usepackage[unicode,hidelinks,bookmarks=false]{hyperref}
\newcommand{\ie}{i.e.\ }
\newcommand{\cf}{cf.\ }
\newcommand{\resp}{respectively\ }
\newcommand{\Subsection}{\subsection}
\providecommand{\nobreakdash}{\nobreak\hbox{-}}
\setlength{\emergencystretch}{2em}
\multlinegap0pt
\usepackage{amssymb}
\usepackage{enumerate}
\usepackage[shortlabels]{enumitem}
\usepackage{algorithm}\usepackage{algpseudocode}
\usepackage[all]{xy}
\usepackage{tikz}
\newtheorem{thm}{Theorem}
\renewcommand{\leq}{\leqslant}
\renewcommand{\phi}{\varphi}
\title{Automorphically Equivalent Elements Of Finite Abelian Groups}
\author{Arjun Agarwal, Rachel Chen, Rohan Garg, Jared Kettinger}
\date{Source: arXiv:2510.06013v2 (CC BY 4.0)}
\begin{document}
\maketitle

\noindent\emph{Source and licence.} Automorphically Equivalent Elements Of Finite Abelian Groups. Version arXiv:2510.06013v2 (CC BY 4.0), distributed by arXiv under the Creative Commons Attribution 4.0 International licence, http://creativecommons.org/licenses/by/4.0/, as recorded in the arXiv licence metadata for this exact version and shown on the abstract page https://arxiv.org/abs/2510.06013v2. Full licence notice: \texttt{licenses/arxiv-2510.06013-CC-BY-4.0.txt}.\par\smallskip

\noindent\textit{Editorial scope.} Counted unit: the article's thm thm (source0.tex lines 272--274). The article's complete original proof of it is delivered with it (source0.tex lines 275--286, one \texttt{proof} environment of 12 lines). No mathematics is written, repaired or re-derived: the statement and the proof are the article's own lines, in the article's own notation, and no retained line is substituted or re-flowed.  No further source-local context is needed: every cross-reference of every retained line resolves inside the two delivered spans.  Nothing was omitted from any retained span, so the package carries no omission record.  No retained line contains a citation, so the wrapper carries no bibliography and no bibliography entry.  No retained line contains a figure environment, an image-inclusion command or drawing code, so no image is delivered and the package declares no asset dependency.  Editorial formatting only, in addition: an AMS \texttt{amsart} preamble that restates the article's own declarations and macros used by the retained lines, namely the environments \texttt{thm} on separate counters, together with the standard packages those lines need.  The retained environments print in this document as Theorem 1 in order of appearance, whereas the article prints the same items with the numbering of its own full document.  The complete native source of the article as distributed in the arXiv e-print for this exact version is bundled unchanged as \texttt{sources/186576/source0.tex}.
\begin{thm}
If $x$ and $y$ are both of maximal order in $G$, they are automorphic images of one another. 
\end{thm}

\begin{proof}
Let $G \cong C_{m_1}\oplus C_{m_2} \oplus \dots \oplus C_{m_k}$ be a finite abelian group written in its invariant factor decomposition. Let $C_{m_r}$, $C_{m_{r+1}}$, $\dots$, $C_{m_k}$ be the invariant factors of maximal order. For an element $x \in G$ to be of maximal order, there exists $i$ where $r \leq i \leq k$ such that the $i$th component of $x$ is a generator of $C_{m_i}$. Define this generator as $x_i$. We construct an automorphism $\phi_{x}$ composing the automorphism switching the $i^\text{th}$ and the $k^\text{th}$ components (which is an automorphism because $C_{m_i}$ and $C_{m_k}$ have the same order, the maximal order) with the automorphism mapping where we map $x$ to the element where all components are kept the same except the $k$-th which becomes $1$. Define $x' = \phi_x(x)$; $x'$ must be of the form $x' = (x'_1, x'_2, \dots, x'_{n-1}, 1)$. We construct another map $\psi_x : G \rightarrow G$ as follows: \begin{align*}
\psi_x : (1, 0, \dots, 0, 0) &\mapsto (1, 0, \dots, 0, 0)\\
\psi_x : (0, 1, \dots, 0, 0) &\mapsto (0, 1, \dots, 0, 0)\\
&\;\;\vdots\\
\psi_x : (0, 0, \dots, 1, 0) &\mapsto (0, 0, \dots, 1, 0)\\
\psi_x : (x'_1, x'_2, \dots, x'_{n-1}, 1) &\mapsto (0, 0, \dots, 0, 1).
\end{align*}
We claim $\psi_x$ is an automorphism. First of all, $\psi_x$ is a homomorphism by construction because the preimages of all the maps we defined is a spanning set of $G$. Furthermore, $\psi_x$ is surjective because the images of all the defined maps form a minimal spanning set of $G$. Since $G$ is finite, injectivity is implied, so $\psi_x$ must be an automorphism. 

Similarly, $\phi_y$ can be defined mapping $y$ to $y'$, and $\psi_y$ can be defined mapping $y'$ to $(0, 0, \dots, 0, 1)$. An automorphism mapping $x$ to $y$ is therefore $\phi_y^{-1}\circ\psi_y^{-1}\circ\psi_x\circ\phi_x$. 
\end{proof}

\end{document}
